\section{Conclusão}

Este trabalho avaliou, por simulação, a atribuição conjunta de solicitações já recebidas a veículos elegíveis, pelo algoritmo Húngaro em lotes de 60 segundos, em comparação com o despacho pelo veículo mais próximo imediato (A) e em lotes (B). Nas duas avaliações reservadas com entradas históricas da Grande Vitória, a atribuição conjunta teve menos km vazio simulado: em relação a A, 5,96\% em setembro de 2026 (15 dias) e 5,04\% em outubro (25 dias); em relação a B, 2,74\% e 1,40\%. Nos 22 cenários sintéticos, que compartilham estrutura e cujas sementes finais foram examinadas antes da adoção do método, a redução mediana foi de 8,4\% em relação a A, evidência exploratória e complementar. Essas são reduções simuladas, de km estimado, e não economia observada.

O atendimento exige a mesma franqueza. Em relação ao despacho em lotes, o critério de atendimento foi preservado nos dois períodos reservados. Em relação ao despacho imediato, foi preservado em setembro (16 atendimentos a mais), mas \textbf{não em outubro}, com 85 atendimentos a menos (14.466 contra 14.551; −0,36 ponto percentual), diferença cujo intervalo no nível diário inclui zero, o que não demonstra equivalência, e que foi reprovada pelo critério fixado antes. Parte da redução de km em outubro pode, portanto, decorrer de menos corridas atendidas; o km por atendido também caiu (−4,5\%), o que sugere efeito além do volume, sem decomposição formal.

As conclusões valem para este simulador, estes dados e estes parâmetros: uma operação, uma região, quatro períodos, distância estimada com fator constante, jornadas reconstruídas e um despacho de referência simplificado, que não reproduz a operação manual. O atraso do lote de 60 s é custo real que o simulador conta na espera. O ótimo por lote não implica ótimo diário. O estudo não demonstra economia na operação real nem generaliza para outras empresas ou cidades; a aplicação prática exigiria validação com a operação, medição do aceite e do cancelamento, distâncias viárias e um piloto controlado. Como trabalho futuro, são naturais a investigação da diferença entre setembro e outubro, a sensibilidade ao tamanho do lote e a comparação com o despacho efetivamente praticado.
